\chapter{Operadores compactos}\label{chap:23}

Un operador acotado controla el tamaño de las imágenes, pero no evita que estas recorran infinitas direcciones separadas. Un operador compacto tiene una propiedad más fuerte: toda sucesión acotada de entradas produce alguna subsucesión convergente de salidas. En ese sentido, recupera parte del comportamiento de la dimensión finita sin exigir que su imagen tenga dimensión finita.

Usaremos la equivalencia entre compacidad y compacidad secuencial en espacios métricos, la acotación total, la completitud y los resultados de dimensión finita del capítulo~\ref{chap:22}. La topología sobre los espacios de operadores será siempre la de la norma de operador, salvo indicación explícita.

\section{Definición y criterios de compacidad}

\begin{definicion}
Un operador lineal $T:V\to W$ es \emph{compacto} si para toda sucesión acotada $(v_n)$ de $V$, la sucesión $(Tv_n)$ tiene una subsucesión convergente en $W$. Denotamos por $\Kcal(V,W)$ el conjunto de estos operadores.
\end{definicion}
\index{operador!compacto}

La frase ``en $W$'' forma parte de la definición. Si $W$ no es completo, una sucesión de imágenes podría ser de Cauchy sin converger en el codominio; eso no basta para la compacidad.

\begin{proposicion}\label{ch23:compacto-acotado}
Todo operador compacto es continuo.
\end{proposicion}
\begin{proof}
Si $T$ no fuera continuo, el criterio del capítulo~\ref{chap:20} permitiría elegir $v_n$ con $\norm{v_n}=1$ y $\norm{Tv_n}>n$. La sucesión de entradas es acotada, pero ninguna subsucesión de las imágenes puede converger: a lo largo de cualquier subsucesión sus normas siguen tendiendo a infinito, mientras que toda sucesión convergente es acotada. Esto contradice la definición.
\end{proof}

\begin{proposicion}\label{ch23:bola-compacta}
Un operador lineal $T:V\to W$ es compacto si y solo si $\cl{T(\overline B(0,1))}$ es compacto en $W$.
\end{proposicion}
\begin{proof}
Suponga primero que $T$ es compacto. Sea $(y_n)$ una sucesión en el cierre de la imagen de la bola. Para cada $n$ elija $v_n$ con $\norm{v_n}\leq1$ y $\norm{Tv_n-y_n}<1/n$. Existe una subsucesión $Tv_{n_j}\to y$ en $W$. Entonces
\[
 \norm{y_{n_j}-y}\leq1/n_j+\norm{Tv_{n_j}-y}\to0.
\]
El límite pertenece al cierre de la imagen porque este es cerrado. Hemos probado que ese cierre es secuencialmente compacto y, por ser métrico, compacto.

Recíprocamente, si el cierre es compacto y $(v_n)$ es acotada, elija $R>0$ con $\norm{v_n}\leq R$ para todo $n$. La sucesión $T(v_n/R)$ pertenece a un conjunto de cierre compacto y tiene una subsucesión convergente. Multiplicar por $R$ da una subsucesión convergente de $(Tv_n)$.
\end{proof}

\begin{proposicion}\label{ch23:totalmente-acotado}
Si $W$ es de Banach, un operador lineal $T:V\to W$ es compacto si y solo si manda conjuntos acotados en conjuntos totalmente acotados.
\end{proposicion}
\begin{proof}
Si $T$ es compacto y $A$ es acotado, existe $R>0$ con $A\subseteq R\overline B(0,1)$. La proposición anterior muestra que $\cl{T(A)}$ está contenido en un compacto, por lo que $T(A)$ es totalmente acotado.

Si $T$ manda conjuntos acotados en totalmente acotados, la imagen de la bola unitaria es totalmente acotada. Su cierre también lo es. Por ser un cerrado del espacio completo $W$, ese cierre es completo. Un espacio métrico completo y totalmente acotado es compacto, por lo que la proposición anterior da la compacidad de $T$. La completitud solo se utilizó en esta segunda dirección.
\end{proof}

\section{Ejemplos y dimensión finita}

\begin{proposicion}\label{ch23:rango-finito}
Todo operador lineal \emph{acotado} de rango finito es compacto. En particular, un operador acotado es compacto cuando su dominio o su codominio tiene dimensión finita.
\end{proposicion}
\begin{proof}
Si $F=T(V)$ es finito dimensional, la imagen de una sucesión acotada es acotada en $F$ porque $T$ es acotado. Una bola cerrada suficientemente grande de $F$ es compacta por Heine--Borel. Existe por ello una subsucesión de imágenes que converge en $F$, y por tanto en $W$.

Si el dominio es finito dimensional, todo operador lineal es continuo y su rango es finito dimensional. Si el codominio es finito dimensional, el rango también lo es, pero la continuidad todavía debe suponerse.
\end{proof}

\begin{observacion}[Corrección de una hipótesis]
La lectura original afirma que el rango finito basta por sí solo. Es falso. En $c_{00}$ con norma uniforme, la funcional
\[
 \varphi(x)=\sum_{n=1}^{\infty}n x_n
\]
está bien definida porque cada $x$ tiene soporte finito, es lineal y tiene rango de dimensión $1$. Sin embargo, $\norm{e_n}_\infty=1$ y $\varphi(e_n)=n$. No es continua ni compacta. La hipótesis ``acotado'' en la proposición anterior corrige exactamente esta dificultad.
\end{observacion}

\begin{ejemplo}[Un operador compacto de rango infinito]\label{ch23:diagonal}
En $\ell^2$, defina $T(x)_k=x_k/k$. La serie que define su norma converge porque
\[
 \sum_{k=1}^{\infty}\left|\frac{x_k}{k}\right|^2\leq\sum_{k=1}^{\infty}|x_k|^2.
\]
Así $T$ lleva $\ell^2$ en sí mismo. Es lineal coordenada a coordenada, $\norm T\leq1$ y, al evaluar en $e_1$, $\norm T=1$.

Para probar su compacidad directamente, sea $(x^{(n)})$ acotada, con $\norm{x^{(n)}}_2\leq R$. Si $R=0$ todas las entradas son cero. Supongamos $R>0$. Para cada coordenada $k$, la sucesión de escalares $x^{(n)}_k/k$ es acotada. Mediante extracciones sucesivas y una diagonal obtenemos una subsucesión $(Tx^{(n_j)})$ cuyas coordenadas convergen todas. La extracción se hace de modo que la subsucesión de la etapa $k+1$ esté contenida en la de la etapa $k$; tomar el término $j$ de la etapa $j$ conserva la convergencia de cada coordenada desde un índice en adelante.

La convergencia de coordenadas todavía no implica convergencia en $\ell^2$. Controlamos las colas uniformemente:
\[
 \sum_{k>N}|(Tx^{(n)})_k|^2
 \leq\frac1{(N+1)^2}\sum_{k>N}|x^{(n)}_k|^2
 \leq\frac{R^2}{(N+1)^2}.
\]
Dado $\eps>0$, elija $N$ con $2R/(N+1)<\eps/2$. Para dos términos de la subsucesión, la norma de la diferencia de sus colas es a lo sumo $2R/(N+1)$. Como las primeras $N$ coordenadas convergen, la norma euclidiana de la diferencia de esas coordenadas es menor que $\eps/2$ cuando ambos índices son grandes. Sumando ambas estimaciones, la subsucesión es de Cauchy en $\ell^2$. La completitud de $\ell^2$ da su convergencia. Por tanto $T$ es compacto. Su rango es infinito dimensional porque contiene $T(e_k)=e_k/k$, una familia linealmente independiente.
\end{ejemplo}

\begin{teorema}\label{ch23:identidad}
La identidad de un espacio normado $V$ es compacta si y solo si $V$ tiene dimensión finita.
\end{teorema}
\begin{proof}
Si la dimensión es finita, la identidad es acotada y de rango finito, luego compacta. Recíprocamente, si es compacta, la proposición~\ref{ch23:bola-compacta} dice que el cierre de $I(\overline B(0,1))=\overline B(0,1)$ es compacto. La bola ya es cerrada, así que es compacta. El teorema~\ref{ch22:caracterizacion} implica dimensión finita.
\end{proof}

\section{Separabilidad y límites en norma}

\begin{teorema}\label{ch23:separable}
Si $T:V\to W$ es compacto, tanto $T(V)$ como $\cl{T(V)}$ son separables. No hace falta que $W$ sea completo.
\end{teorema}
\begin{proof}
Por la proposición~\ref{ch23:bola-compacta}, los conjuntos
\[
 K_m=\cl{T(\overline B(0,m))}\qquad(m\in\N)
\]
son compactos. Todo compacto métrico es separable: para cada $n$ se puede elegir una red finita de radio $1/n$, y la unión de todas esas redes es numerable y densa. Por tanto cada subespacio métrico $T(\overline B(0,m))$ es separable; el argumento de base numerable del ejercicio 4 del capítulo~\ref{chap:19} justifica esta herencia a subconjuntos.

Elija un conjunto numerable $D_m$ denso en $T(\overline B(0,m))$. Como
\[
 T(V)=\bigcup_{m=1}^{\infty}T(\overline B(0,m)),
\]
la unión $D=\bigcup_mD_m$ es numerable y densa en $T(V)$. En efecto, cada punto de $T(V)$ pertenece a alguna de esas imágenes y puede aproximarse con puntos de su $D_m$. Tomando cierres, $\cl D=\cl{T(V)}$, de modo que el mismo conjunto es denso en el cierre.
\end{proof}

\begin{proposicion}
El conjunto $\Kcal(V,W)$ es un subespacio de $\Bcal(V,W)$.
\end{proposicion}
\begin{proof}
El operador cero es compacto. Si $S,T$ son compactos y $(v_n)$ es acotada, extraiga primero una subsucesión en la que $Sv_n$ converja. De esa subsucesión extraiga otra en la que $Tv_n$ converja. En la segunda, ambas imágenes convergen, y por la continuidad de las operaciones también converge $\alpha Sv_n+\beta Tv_n$. Esto vale para cualesquiera escalares $\alpha,\beta$, por lo que $\alpha S+\beta T$ es compacto.
\end{proof}

\begin{teorema}\label{ch23:cierre}
Si $W$ es de Banach y $T_n\in\Kcal(V,W)$ converge en norma de operador a $T\in\Bcal(V,W)$, entonces $T$ es compacto. En consecuencia, $\Kcal(V,W)$ es un subespacio cerrado y completo de $\Bcal(V,W)$.
\end{teorema}
\begin{proof}
Demostraremos que $T(\overline B(0,1))$ es totalmente acotado. Dado $\eps>0$, elija $n$ con $\norm{T-T_n}<\eps/3$. Como $T_n$ es compacto, su imagen de la bola unitaria es totalmente acotada. Existen $y_1,\ldots,y_k\in W$ tales que cada $T_nv$, con $\norm v\leq1$, está a distancia menor que $\eps/3$ de alguno de ellos. Para ese centro,
\[
 \norm{Tv-y_j}\leq\norm{(T-T_n)v}+\norm{T_nv-y_j}<2\eps/3<\eps.
\]
La familia finita de bolas $B(y_j,\eps)$ cubre la imagen por $T$. Reescalar muestra lo mismo para la imagen de cualquier bola, y todo conjunto acotado está contenido en una bola. Por la proposición~\ref{ch23:totalmente-acotado}, $T$ es compacto.

La cerradura se sigue del criterio de sucesiones en el espacio normado de operadores. Como $\Bcal(V,W)$ es completo por el teorema~\ref{ch20:banach-operadores}, su subespacio cerrado $\Kcal(V,W)$ también lo es.
\end{proof}

\begin{observacion}
El operador del ejemplo~\ref{ch23:diagonal} también se aproxima en norma por sus truncaciones $T_N(x)=(x_1,x_2/2,\ldots,x_N/N,0,\ldots)$. Estas son de rango finito y
\[
 \norm{T-T_N}=\frac1{N+1}.
\]
La cota superior se obtiene de la estimación de colas, y el vector $e_{N+1}$ da la cota inferior. El teorema de cerradura proporciona así una segunda prueba de compacidad. No se está afirmando que todo operador compacto entre espacios de Banach arbitrarios admita aproximación por operadores de rango finito; el resultado probado va en la dirección de la aproximación hacia la compacidad.
\end{observacion}

\begin{ejemplo}[Por qué importa el codominio completo]
Considere $V=W=c_{00}$ con norma uniforme y $Tx=(x_k/k)_k$. Sus truncaciones $T_N$ tienen rango finito, son acotadas y compactas. Además $\norm{T-T_N}=1/(N+1)\to0$. Sin embargo, $T$ no es compacto con valores en $c_{00}$. Las entradas $x^{(m)}=(1,\ldots,1,0,\ldots)$, con $m$ unos, tienen norma $1$. Sus imágenes convergen uniformemente en $c_0$ a $(1,1/2,1/3,\ldots)$, que no tiene soporte finito. Ninguna subsucesión puede converger en $c_{00}$, porque cualquier límite tendría necesariamente esas mismas coordenadas. Así la hipótesis de completitud en el teorema~\ref{ch23:cierre} no es decorativa.
\end{ejemplo}

\section{Ejercicios resueltos}

\begin{ejercicio}[Lectura 23, ejercicio 1]
Si $f:V\to W$ es uniformemente continua en todo el espacio normado $V$ y $A\subseteq V$ es acotado, demuestre que $f(A)$ es acotado.
\end{ejercicio}
\begin{solucion}
Por continuidad uniforme existe $\delta>0$ tal que $\norm{x-y}<\delta$ implica $\norm{f(x)-f(y)}<1$. Elija $R>0$ con $\norm a\leq R$ para $a\in A$ y un entero $N$ con $R/N<\delta$. Para cada $a\in A$, divida el segmento del origen a $a$ en los puntos $a_j=ja/N$, $0\leq j\leq N$. Cada salto tiene norma $\norm{a_j-a_{j-1}}=\norm a/N<\delta$. Por tanto
\[
 \norm{f(a)-f(0)}\leq\sum_{j=1}^N\norm{f(a_j)-f(a_{j-1})}<N.
\]
Así $\norm{f(a)}\leq\norm{f(0)}+N$ para todos los $a\in A$. Es una sola cota porque $N$ se eligió a partir de $R$, no de cada punto. La prueba utiliza que el segmento completo está en el dominio; para funciones definidas solamente en un conjunto métrico arbitrario la afirmación no tendría este argumento.
\end{solucion}

\begin{ejercicio}[Lectura 23, ejercicio 2]
Si $W$ es de Banach, pruebe que un operador lineal $T:V\to W$ es compacto si y solo si manda conjuntos acotados en conjuntos totalmente acotados.
\end{ejercicio}
\begin{solucion}
La proposición~\ref{ch23:totalmente-acotado} ya da la demostración completa. Podemos expresar el núcleo del argumento mediante sucesiones. Si la imagen de algún acotado $A$ no fuera totalmente acotada, existiría $\eps>0$ y podríamos elegir inductivamente $y_n=Tx_n$ con $x_n\in A$ y $\norm{y_n-y_m}\geq\eps$ para $n\ne m$: cada nuevo punto se elige fuera de las bolas de radio $\eps$ centradas en los anteriores. Ninguna subsucesión de $(y_n)$ sería de Cauchy, contradiciendo compacidad. En sentido inverso, una sucesión en un conjunto totalmente acotado tiene una subsucesión de Cauchy, obtenida escogiendo sucesivamente bolas de radios tendiendo a cero que contengan infinitos términos. La completitud de $W$ transforma esa subsucesión de Cauchy en una convergente en $W$.
\end{solucion}

\begin{ejercicio}[Lectura 23, ejercicio 3]
Sean $S:V\to W$ y $T:W\to X$ operadores continuos. Si uno de ellos es compacto, demuestre que $TS$ es compacto.
\end{ejercicio}
\begin{solucion}
Sea $(v_n)$ acotada en $V$. Si $S$ es compacto, existe una subsucesión $Sv_{n_j}\to w$ en $W$. La continuidad de $T$ da $TSv_{n_j}\to Tw$ en $X$.

Si $T$ es compacto, la continuidad lineal de $S$ da $\norm{Sv_n}\leq\norm S\sup_m\norm{v_m}$, de modo que $(Sv_n)$ es acotada. La compacidad de $T$ aplicada a esta sucesión produce una subsucesión convergente de $(TSv_n)$. En ambos casos se verifica la definición. No se requiere que los espacios sean completos.
\end{solucion}

\begin{ejercicio}[Lectura 23, ejercicio 4]
Sea $T:V\to V$ compacto y $V$ infinito dimensional. Pruebe que no existe un operador continuo $S$ con $ST=TS=I$.
\end{ejercicio}
\begin{solucion}
Si existiera, el ejercicio anterior haría compacto a $ST=I$. Por el teorema~\ref{ch23:identidad}, eso obligaría a que $V$ tuviera dimensión finita, contradicción. De hecho, basta la igualdad $ST=I$ para llegar a la contradicción. El resultado se refiere a una inversa continua: una biyección algebraica no garantiza automáticamente una inversa continua en espacios normados arbitrarios.
\end{solucion}

\section{Perspectiva}
Los operadores compactos reemplazan la acotación por una propiedad que permite extraer límites. Se comportan bien bajo composición con operadores acotados y, cuando el codominio es completo, bajo límites en norma. La identidad muestra el alcance de la noción: ser compacto es una restricción real en dimensión infinita, aunque muchos operadores que atenúan coordenadas sí la satisfacen.
